% Standalone results draft for the JMD extension.
% Compile from either the repository root or experiment/analysis.
% Existing AV figures are referenced in place; organizational assets are staged
% in paper_figures/idetc. No new organizational interventions are reported.
% Main manuscript text ends before the author-only notes.
\documentclass[11pt]{article}
\usepackage[margin=0.9in]{geometry}
\usepackage[T1]{fontenc}
\usepackage{lmodern}
\usepackage{graphicx}
\usepackage{amsmath}
\usepackage{booktabs}
\usepackage{tabularx}
\usepackage{caption}
\usepackage{placeins}
\usepackage[hidelinks]{hyperref}
\captionsetup{font=small,labelfont=bf}
\setlength{\parskip}{0.35em}
\setlength{\parindent}{0pt}
\setlength{\emergencystretch}{2em}

\newcommand{\AnalysisRoot}{.}
\IfFileExists{experiment/analysis/track_social_effects/results/figures/complexity_effects.pdf}
  {\renewcommand{\AnalysisRoot}{experiment/analysis}}{}
\newcommand{\ResultGraphic}[2][0.97]{%
  \includegraphics[width=#1\linewidth,height=0.53\textheight,keepaspectratio]
    {\AnalysisRoot/#2}}
\newif\ifAuthorNotes
\AuthorNotestrue % Change to \AuthorNotesfalse when preparing manuscript text.

\title{When social relationships change technology use\\
  \large Results and cross-case interpretation for the JMD extension}
\author{}
\date{}

\begin{document}
\maketitle
\vspace{-2em}

\section{Results}
\label{sec:results}

The analysis addresses the question: \emph{How much do social relationships
shape decisions and technology use in sociotechnical systems, and under what
conditions do they matter most?} We examine this question through a vehicle
model with participant-specific individual responses and an organizational
model developed for AI adoption in engineering teams~\cite{noorghasemi2026}.
The vehicle comparisons quantify the change produced by adding modeled peer
influence. The organizational results describe collective outcomes within a
socially coupled system. These provide complementary evidence, with different
outcomes and levels of individual detail.

\subsection{Measuring the social effect on driver decisions}
\label{sec:measure}

The common vehicle outcome is the driver's executed choice between automation
and manual driving. Each social run is compared with a solo run using the same
driver model, initial individual state and track. For a drive of $T$ steps,
the social effect is
\begin{equation}
  E_m = \frac{100}{T}\sum_{t=1}^{T}
    \mathbf{1}\!\left[d_m(t)\ne d_{\mathrm{solo}}(t)\right],
  \label{eq:social-effect}
\end{equation}
where $m$ is the influence mechanism and $d(t)$ is the binary driver decision.
Thus, an effect of 18.3\% means that approximately 18 of every 100 decisions
differ from the matched solo decisions. Changes in either direction count.
This measure describes a change in predicted behavior; it does not measure an
improvement in safety or decision quality.

Decision-level influence combines occupants' expressed choices after their
individual decision rules have been applied. Cognition-level influence combines
the continuous states entering those rules: automation trust, perceived risk
and workload. In the cognition-level model, the socially altered states also
carry into later steps. The comparison therefore includes both the location of
influence and its accumulated consequences. These mechanisms do not measure
the duration or depth of an actual conversation.

Figure~\ref{fig:examples} illustrates why the social layer cannot be assigned
one effect size for every drive. In one selected case, both social models differ
from solo at only one of the 100 steps. In another, decision-level and
cognition-level influence change 49 and 47 decisions, respectively. The cases
use the same driver and track but different passenger and starting-state
configurations. They illustrate variation; the following aggregate comparisons
establish how the effect changes across conditions.

\begin{figure}[htbp]
  \centering
  \ResultGraphic{example_cases/results/figures/example_cases_simple.pdf}
  \caption{Driver automation choices in three matched model variants for two
  selected cases. Blue denotes decision-level influence, orange denotes
  cognition-level influence, and the solo trajectory provides the reference.
  Gray shading identifies high task-complexity steps. Small vertical offsets
  distinguish coincident lines. The two panels are illustrative cases, not
  estimates of how frequently these patterns occur.}
  \label{fig:examples}
\end{figure}
\FloatBarrier

\subsection{Track complexity changes the magnitude of the social effect}
\label{sec:complexity}

The mean social effect decreases across the Low, Medium and High track-complexity
conditions (Fig.~\ref{fig:complexity}). Decision-level influence changes 18.3\%,
10.5\% and 6.1\% of driver decisions, respectively. Cognition-level influence
changes 27.5\%, 21.9\% and 18.7\%. Cognition-level influence therefore produces
larger departures from solo at each complexity level, and its effect remains
substantial in the High condition.

At High complexity, cognition-level influence changes approximately three times
as many decisions as decision-level influence. The 18.7\% cognition-level effect
at High complexity is also roughly comparable to the 18.3\% decision-level
effect at Low complexity. High complexity therefore does not by itself imply
a small modeled social effect. The near equality of these two values is
descriptive; it does not establish statistical equivalence.

These averages use the same cars and starting configurations across tracks,
give equal weight to drivers and passenger counts, and average over the
persistence settings. Persistence is treated here as a secondary scenario
variable, rather than assumed to have no effect. The nine fixed trajectories
span three complexity levels and three persistence settings; repeated cars on
these tracks are not independent replications of track generation.

\begin{figure}[htbp]
  \centering
  \ResultGraphic[0.679]{track_social_effects/results/figures/complexity_effects.pdf}
  \caption{Mean percentage of driver decisions different from matched solo
  driving at Low, Medium and High track complexity. Blue denotes decision-level
  influence and orange denotes cognition-level influence. The analysis uses
  514,800 matched base configurations on nine tracks, averages over persistence
  settings, and balances drivers and passenger counts. Percentages describe
  changed choices, not gains in performance.}
  \label{fig:complexity}
\end{figure}
\FloatBarrier

\subsection{The social effect develops differently across steps and complexity}
\label{sec:time}

Whole-drive averages conceal different temporal patterns
(Fig.~\ref{fig:timing}). On High-complexity tracks, the decision-level effect
falls after the opening steps, whereas the cognition-level effect remains
larger later in the drive. On Low-complexity tracks, social effects persist;
the cognition-level effect becomes larger after the opening period. Social
influence therefore does not uniformly fade as the drive proceeds.

A separate early/later aggregation of the full comparison grid confirms this
pattern. At High complexity, decision-level differences fall from 15.7\% in
steps 1--20 to 3.9\% in steps 21--100; cognition-level differences fall from
22.3\% to 19.2\%. At Low complexity, decision-level differences are nearly
unchanged (18.4\% to 18.8\%), while cognition-level differences increase from
24.7\% to 29.7\%. Each percentage is a rate per step, so the longer later window
does not itself increase the rate. These full-grid values use pooled case
weights; the ten-step curves use a separate sample, as specified in
Appendix~\ref{app:samples}.

Together, the complexity and time results show that a single average social
effect would miss two relevant conditions: the external task context and where
peer influence enters the individual's response. The larger later effects under
cognition-level influence are consistent with its state carryover, but the
present comparison does not isolate carryover from current-step influence.

\begin{figure}[htbp]
  \centering
  \ResultGraphic{influence_timing/results/figures/step_blocks_by_complexity.pdf}
  \caption{Driver decisions different from solo within successive ten-step
  blocks, grouped by influence mechanism and track complexity. Points
  represent steps 1--10, 11--20, through 91--100 and are placed at block
  midpoints. Colors identify track complexity; solid lines show decision-level
  influence and dashed lines show cognition-level influence. Curves use a fixed sample of 46,800 car--track cases
  (200 per driver and track), with pooled case weights. Lines connect block
  averages.}
  \label{fig:timing}
\end{figure}
\FloatBarrier

\subsection{Social feedback changes confidence in the driver's own judgment}
\label{sec:confidence}

The social layer changes self-trust as well as the current automation choice.
Trajectories grouped by initial confidence are shown in
Fig.~\ref{fig:starting-confidence}. From a common starting mean of 0.515, mean final
self-trust reaches 0.798 with decision-level influence and 0.673 with
cognition-level influence. Solo self-trust remains constant under the model's
update rule. Self-trust is confidence in one's own judgment, distinct from
trust in automation. Here it is updated from passengers' trust in the driver,
not from measured driving success or decision correctness.

This result connects confidence to the paper's central question: the social
layer alters a variable that helps determine how much weight the driver gives
their own judgment relative to passengers' judgments. It changes both choices
and the conditions for later influence. Higher confidence in these simulations
does not establish better-calibrated confidence.

Initial-confidence categories reveal narrowing differences within each
influence model (Fig.~\ref{fig:starting-confidence}; category definitions in
Appendix~\ref{app:confidence}). The Low, Medium and High groups end with
closely spaced means in decision mode (0.796, 0.796 and 0.803). The corresponding
cognition-mode means are 0.677, 0.651 and 0.685. These are group averages rather
than evidence that every driver's confidence converges. Final confidence is
also higher on higher-complexity tracks in both social models, although those
tracks have fewer decisions changed from solo. These associations do not
identify increased confidence as the cause of the smaller social effect.

\begin{figure}[htbp]
  \centering
  \ResultGraphic{driver_confidence/results/figures/starting_confidence.pdf}
  \caption{Mean self-trust grouped by Low, Medium and High initial driver
  confidence. Differences between group means narrow within each influence
  model. Step 0 is the initial state, followed by post-update states through
  step 100. The same 780 sampled base configurations are repeated on nine
  tracks, with equal driver and passenger-count weights. Groups describe the
  sampled configurations; other initial conditions are not held fixed between
  categories. These curves do not establish that every driver reaches a common
  confidence value.}
  \label{fig:starting-confidence}
\end{figure}
\FloatBarrier

\subsection{Weak initial relationships can become consequential through learning}
\label{sec:weak-ties}

A controlled experiment tests whether weak initial passenger ties keep driver
behavior close to solo driving (Fig.~\ref{fig:weak-ties}). Interpersonal trust
is reduced by a factor of 1,024 while initial self-trust is retained. Occupants,
initial individual states and tracks remain matched. Both social mechanisms
then reproduce the solo choice in the first step.

With trust learning enabled, differences reach 9.1\% of decisions under
decision-level influence and 21.9\% under cognition-level influence in steps
81--100. Holding trust fixed gives corresponding rates of 0.0\% and 7.5\%.
An initially small social effect therefore does not guarantee a small later
effect. The remaining cognition-level differences with fixed trust show that
small, nonzero social weights can still change later decisions when cognitive
states are carried forward.

The fixed condition freezes both interpersonal trust and self-trust updates;
the comparison estimates their combined contribution. Learning also does not
always amplify influence: with the original trust values, the late
decision-level effect is 9.1\% with learning and 13.5\% with fixed trust.
The finding is specifically that learning can make initially weak relationships
consequential, rather than that every learning process increases the social
effect. This targeted experiment uses 130 starting configurations on the nine
tracks and is separate from the full starting-state sweep.

\begin{figure}[htbp]
  \centering
  \ResultGraphic{trust_learning/results/figures/weak_ties.pdf}
  \caption{Social effects with original relationships and with weak initial
  interpersonal ties. Weak ties retain initial self-trust and reduce every
  interpersonal trust value by a factor of 1,024. Learning allows both
  interpersonal trust and self-trust to change; the fixed condition freezes
  both. Each curve shows the percentage of decisions different from matched
  solo driving at each step, across 1,170 targeted car--track cases. Weak ties
  preserve network links and small nonzero weights.}
  \label{fig:weak-ties}
\end{figure}
\FloatBarrier

\subsection{Organizational initial attitudes change the performance threshold}
\label{sec:organization}

The organizational analysis examines a different consequence of socially
coupled judgments: how technical performance and initial attitudes relate to
a collective outcome. Engineers revise a scalar attitude through peer
influence and satisfactory or unsatisfactory AI outcomes. The AI success
probability is the Bernoulli probability of a satisfactory outcome; it is not
a measured accuracy for a specific deployed AI tool.

We retain the conference framework's threshold analysis, with its outcome
stated more precisely. The reported sum of quiet and loud adopter fractions,
$Q+L$, counts agents with nonnegative AI attitudes. We call this the
\emph{favorable-attitude fraction}. It differs from actual AI use because
some agents with negative attitudes still meet the model's use threshold.
Consequently, the threshold in Fig.~\ref{fig:org-threshold} is the AI success
probability needed for a surrogate-predicted steady-state favorable-attitude
fraction of at least 50\%. The surrogate predicts rates for the reduced
stochastic model, whose stationary distribution supplies this outcome. The
threshold is selected on a probability grid with increments of 0.001.

With a positive mean initial attitude of $+0.25$, this threshold ranges from
approximately 0.67 to 0.71 across the plotted team configurations. With a
negative mean initial attitude of $-0.25$, it ranges from 0.71 to 0.91. Structural
variation is therefore associated with a wider range of thresholds when the
organization starts negatively. The map does not show a uniformly monotonic
effect of increasing either team size or team count.

This result adds a collective design consequence to the individual vehicle
comparisons: the technical performance needed to reach a specified attitude
target depends on the initial social setting and the organizational
configuration. All plotted organizational conditions include the social model,
however. The threshold map does not isolate how much of that dependence is
caused by peer influence relative to a matched individual-only organization.

\begin{figure}[htbp]
  \centering
  \ResultGraphic{paper_figures/idetc/structure_min_accuracy_clarified.pdf}
  \caption{Minimum AI success probability for a surrogate-predicted steady-state
  favorable-attitude fraction of at least 50\%, across team sizes and numbers
  of teams. Mean initial attitudes are $-0.25$ (left) and $+0.25$ (right).
  Favorable attitudes are the original organizational model's $Q+L$ outcome.
  This figure retains the threshold analysis developed
  for the conference framework~\cite{noorghasemi2026}, with clarified outcome
  labels. It compares configurations within the social model, rather than
  social and individual-only variants.}
  \label{fig:org-threshold}
\end{figure}

The organizational composition result provides a secondary distinction
(Appendix~\ref{app:organization}): having a favorable attitude and having a
high modeled influence score are different outcomes. Their relative shares
vary with AI success probability and initial sentiment. We retain this
distinction without interpreting a high influence score as observed advocacy
or as proof that these agents caused adoption by others.
\FloatBarrier

\section{Cross-case interpretation}
\label{sec:interpretation}

The two applications are comparable through their modeled processes rather
than through identical tasks or timescales (Table~\ref{tab:comparison}). In
both, people form judgments about a technology, weight peers' judgments through
trust, make repeated choices, and update relationships. The vehicle case
resolves individual judgment into participant-fitted states and decision
rules; the organizational case represents attitude more coarsely and traces
collective outcomes. Vehicle observations inform individual response models,
while the social interactions in both applications remain simulated mechanisms.
Passenger agents reuse the response models fitted to drivers, making their
responses as passengers an additional modeling assumption.
The organizational simulation is not empirically validated by the vehicle data.

\begin{table}[htbp]
  \centering
  \small
  \caption{Functional correspondences and differences between the cases.}
  \label{tab:comparison}
  \begin{tabularx}{\linewidth}{@{}lXX@{}}
    \toprule
    Dimension & Vehicle & Engineering organization \\
    \midrule
    Individual judgment & Automation trust, risk and workload & Scalar attitude toward AI \\
    Peer influence & Influence on states or expressed choices & Influence on attitude \\
    Choice authority & One driver executes the shared choice & Engineers make distributed use choices \\
    Structure & Driver/passenger role asymmetry & Agents, teams and organizational aggregation \\
    Technical context & Track/task complexity & AI satisfactory-outcome probability \\
    Development & Repeated driving steps & Repeated AI interactions \\
    \bottomrule
  \end{tabularx}
\end{table}

The vehicle comparisons establish that the contribution of social
relationships is conditional: its magnitude depends on influence location,
track complexity and interaction history. The organizational results identify
a collective consequence to examine with the same approach: the performance
threshold varies with initial attitudes and team configuration within the
social model. A matched organizational counterfactual must establish how much
peer influence contributes to that variation. This connects the research
question across scales without claiming that both cases already demonstrate
the same causal pattern.

Influence on underlying judgment is a plausible organizational mechanism.
Research on engineering teams connects coordination with changes in shared
understanding~\cite{avnet2016}, and experimental work shows how information
exchange and integration can alter group decisions~\cite{vanginkel2008}.
This supports considering more than the pooling of expressed preferences.
It does not assign engineering organizations to the vehicle cognition-level
curve: neither model measures discussion depth, and engineering task
complexity has not been calibrated to the track-complexity categories.
The vehicle percentages therefore remain vehicle results.

The weak-ties experiment motivates a specific organizational expectation:
initially weak relationships may become consequential if trust adapts during
repeated interaction. This is a temporal relationship mechanism demonstrated
in the vehicle simulations, and an organizational hypothesis until a matched
fixed/adaptive comparison is performed. The organization also contains a
distinct feedback path: an attitude change can alter AI use, which changes
whether an engineer receives a subsequent AI outcome. That experience may
reinforce or reverse the socially induced attitude. Testing this path would
connect the individual mechanism to collective adoption dynamics without
assuming identical behavior across the applications.

For design, the vehicle results show why an individual-only model can miss
changes in choices and in later susceptibility to influence. The organizational
thresholds show why specifying technical performance alone may be insufficient
to predict a collective attitude target. The quantitative evidence for the
size of the social contribution currently comes from matched vehicle
comparisons; the organizational results establish a consequence to test with
the same kind of counterfactual comparison.

The scope of the evidence remains the implemented models, fitted individuals
and sampled scenarios. The 26 participant-specific vehicle models and nine
fixed tracks are not enlarged empirical samples by repeated simulation.
Differences between cognitive and decision influence also combine influence
location with state persistence, and the learning intervention combines two
trust-update rules. Independent scenario realizations and separate mechanism
interventions would strengthen broader claims. No physical duration is assigned
to a step in either model, so the temporal comparison concerns repeated
interactions rather than equal elapsed time.
\FloatBarrier

\appendix
\section{Supporting confidence results}
\label{app:confidence}

Three initial-confidence categories keep the descriptive comparison readable:
Low ($c_0<1/3$), Medium ($1/3\leq c_0<2/3$) and High ($c_0\geq2/3$).
These are display categories on the model scale, not validated psychological
thresholds. Other starting conditions differ between categories; this grouping
is not an intervention that changes only confidence.

Figure~\ref{fig:inbound-trust} shows the other side of the confidence feedback:
the mean trust passengers place in the driver. Cars are grouped by their
initial mean passenger trust, rather than by the driver's initial confidence.
Mean passenger trust starts at 0.188, 0.497 and 0.810 in the Low, Medium and
High groups. It ends at 0.778--0.789 with decision-level influence and
0.657--0.668 with cognition-level influence. Differences between initial
passenger-trust group means therefore narrow within each model, as the
initial-confidence group means do in Fig.~\ref{fig:starting-confidence}.
Self-trust moves toward the previous passenger appraisal under the model's
update rule; this relationship is not an independent empirical discovery.

\begin{figure}[htbp]
  \centering
  \ResultGraphic{driver_confidence/results/figures/inbound_trust_by_starting_trust.pdf}
  \caption{Mean simulated passenger-to-driver trust over steps, grouped by
  Low, Medium and High initial mean passenger trust. Category boundaries are
  1/3 and 2/3 on the model scale. Each car contributes its
  mean raw inbound trust from passengers, excluding the driver's self-trust;
  car means then use the same driver and passenger-count weighting as the
  confidence figure. Step 0 is the initial state, followed by post-update
  states through step 100. The quantity is raw trust on the 0--1 model scale,
  rather than summed trust or normalized influence. Other starting conditions
  are not held fixed between groups. Solo has no passenger appraisal.}
  \label{fig:inbound-trust}
\end{figure}
\FloatBarrier

Final mean self-trust is 0.611, 0.851 and 0.933 at Low, Medium and High
complexity under decision-level influence, and 0.562, 0.701 and 0.754 under
cognition-level influence. Figure~\ref{fig:confidence-complexity} concerns
confidence adaptation, whereas Fig.~\ref{fig:complexity} concerns differences
in executed choices. Their opposing associations with complexity do not by
themselves establish a causal link between confidence and the social effect.

\begin{figure}[htbp]
  \centering
  \ResultGraphic{driver_confidence/results/figures/confidence_by_complexity.pdf}
  \caption{Mean driver self-trust by track complexity in the two social
  models. Cars and initial states are matched across tracks, and each
  complexity category averages over persistence settings. Higher final
  confidence on more complex tracks is an association under the specified
  social-feedback rule.}
  \label{fig:confidence-complexity}
\end{figure}
\FloatBarrier

\section{Supporting organizational results and reduction fidelity}
\label{app:organization}

The original quiet/loud classification separates favorable agents by a modeled
influence score. For agent $i$, the score is its attitude multiplied by the
sum of incoming raw trust values, $\ell_i=o_i\sum_j w_{ji}$. Nonnegative-attitude
agents below the score threshold of 0.5 are quiet ($Q$); those at or above it
are loud ($L$). The score is a model proxy, affected by attitude, trust and
the number of incoming connections. It is not a recorded frequency of speaking
or a measured advocacy behavior.

Figure~\ref{fig:org-composition} retains this secondary conference analysis.
Quiet favorable agents are prominent at moderate AI success probabilities in
several positive and neutral conditions; higher-score favorable agents become
more prominent toward the high-success end of the plotted range. The curves
show population fractions, rather than fractions conditional on being
favorable. Their composition does not identify which agents caused others to
change attitude.

\begin{figure}[htbp]
  \centering
  \ResultGraphic{paper_figures/idetc/adopter_composition.pdf}
  \caption{Surrogate-predicted steady-state population fractions in the quiet ($Q$) and loud
  ($L$) favorable-attitude classes, and their sum, across AI success
  probabilities. Rows show positive ($+0.3$, top), neutral ($0$, middle), and
  negative ($-0.3$, bottom) mean initial attitudes. Columns show the plotted team
  configurations. The original labels ``accuracy'' and ``adoption'' correspond
  to satisfactory-outcome probability and favorable-attitude fraction here.
  $Q$ and $L$ denote lower and higher modeled influence scores, not observed
  passive use and advocacy. Asset retained from the conference framework.}
  \label{fig:org-composition}
\end{figure}
\FloatBarrier

\clearpage
\subsection*{Reduction fidelity}

The organizational reduction provides computational support for exploring
these collective outcomes. The conference analysis reports a correlation of
0.994 between final favorable-attitude fractions in the generating agent-based
model and its fitted reduced model (Fig.~\ref{fig:reduction}). Median trajectory
RMSE over all four state fractions and recorded steps is approximately 0.0209,
with an interquartile range of 0.0108--0.0436.
These checks establish internal approximation fidelity to the generating
simulations. They do not directly assess surrogate-predicted stationary
fractions, validate real organizational behavior, or show that the original
reduction preserves effects of social interventions absent from its fitting data.

\begin{figure}[htbp]
  \centering
  \ResultGraphic{paper_figures/idetc/reduction_parity.pdf}
  \caption{Final favorable-attitude fraction in the organizational agent-based
  model and its fitted CTMC/ODE reduction. The original asset labels this
  quantity as adoption. Agreement is a check of reduction fidelity to the
  generating model, rather than empirical validation of AI adoption. Retained
  from the conference framework~\cite{noorghasemi2026}.}
  \label{fig:reduction}
\end{figure}
\FloatBarrier

\section{Comparison populations and aggregation}
\label{app:samples}

The analyses use different samples and weights (Table~\ref{tab:samples}).
Their percentages answer related questions but should not be combined as
though they came from one common population. All vehicle comparisons use the
same nine fixed trajectories and 100 decision steps per drive. Here a base
configuration specifies a driver, passengers and initial model conditions.

\begin{table}[htbp]
  \centering
  \small
  \caption{Vehicle analysis populations and weighting.}
  \label{tab:samples}
  \begin{tabularx}{\linewidth}{@{}l r X@{}}
    \toprule
    Analysis & Car--track cases & Aggregation \\
    \midrule
    Whole-drive complexity & 4,633,200 & 514,800 base configurations; equal driver and passenger-count weights; persistence averaged \\
    Early/later timing & 4,633,200 & Full comparison grid; pooled case weights; steps 1--20 and 21--100 \\
    Ten-step timing & 46,800 & Fixed sample of 200 cases per driver and track; pooled case weights \\
    Confidence & 7,020 & 780 matched base configurations; equal driver and passenger-count weights \\
    Trust-learning intervention & 1,170 & 130 base configurations; one per driver and passenger count; targeted starting states \\
    \bottomrule
  \end{tabularx}
\end{table}
\FloatBarrier

\begin{thebibliography}{9}
\bibitem{noorghasemi2026}
Noorghasemi, A., and McComb, C., 2026,
``Designing for AI Adoption: A Multi-Scale Modeling Framework for Engineering
Organizations,'' \emph{Proceedings of the ASME International Design Engineering
Technical Conferences and Computers and Information in Engineering Conference},
Houston, TX, Paper No. DETC2026-194708.

\bibitem{avnet2016}
Avnet, M. S., 2016, ``A Network-Based Analysis of Team Coordination and Shared
Cognition in Systems Engineering,'' \emph{Systems Engineering}, \textbf{19}(5),
pp. 395--408. \href{https://doi.org/10.1002/sys.21341}{doi:10.1002/sys.21341}.

\bibitem{vanginkel2008}
van Ginkel, W. P., and van Knippenberg, D., 2008,
``Group Information Elaboration and Group Decision Making: The Role of Shared
Task Representations,'' \emph{Organizational Behavior and Human Decision
Processes}, \textbf{105}(1), pp. 82--97.
\href{https://doi.org/10.1016/j.obhdp.2007.08.005}{doi:10.1016/j.obhdp.2007.08.005}.
\end{thebibliography}

\ifAuthorNotes
\clearpage
\section*{Author notes: selection and checks before submission}
\addcontentsline{toc}{section}{Author notes}
\textbf{These notes are outside the manuscript text.} Set
\verb|\AuthorNotesfalse| in the preamble to omit them from the compiled draft.

\textbf{IDETC selection.} Keep the minimum-performance threshold map in the
main results because it supplies a collective consequence and connects social
initial conditions to a technical requirement. Keep quiet/loud composition
as a supporting result in the appendix unless the journal paper explicitly
tests how influence is distributed. Move agent-based/reduced-model parity and
trajectory agreement to methods or supplementary material. Surrogate degree
selection is a methodological detail; it is omitted here. Reduction accuracy
should not be presented as a primary finding about social importance.

\textbf{Organizational counterfactual still needed.} The present organizational
section retains existing socially coupled results; it does not report a new
individual-only comparison. A matched no-peer/fixed-trust/adaptive-trust
experiment would let the journal extension quantify the organizational social
contribution directly. Keep potential AI outcomes matched by agent and step;
a shared random seed alone may not match experiences when use choices differ.
Report actual use and favorable attitudes separately.

\textbf{Implementation and figure provenance require confirmation.} The current
organizational source has commented-out writeback of the inter-team opinion
update, and its update order differs from the conference conceptual diagram.
Identify the exact source version that generated the retained conference
results before claiming demonstrated cross-team behavioral propagation.
Do not infer managerial authority from team aggregation alone. Copied figure
origins and hashes are recorded in \verb|paper_figures/idetc/sources.json|;
the relabeled threshold figure has a separate reproduction script and input
provenance. The retained grid values should be reconciled with the submission's
source archive.

\textbf{Terminology requiring alignment with methods.} State the $Q+L$ outcome
as favorable attitudes, define the quiet/loud score, and distinguish it from
actual technology use. Describe cognition-level influence as influence on
judgment, without asserting measured deep discussion. Do not map engineering
tasks onto the AV complexity axis or project AV percentages into organizations.
The fixed-trust experiment freezes two learning rules; the cognition comparison
also changes state carryover. Separate interventions would identify these
contributions more precisely.

\textbf{Figure selection.} The main sequence is examples, complexity, timing,
confidence grouped by initial confidence, weak initial ties, and the
organizational threshold. Passenger-to-driver trust grouped by its initial
level, confidence by complexity, organizational composition and reduction
parity are supporting figures. Similarity categories, baseline regimes,
pressure diagnostics and the old network averages are omitted because they
do not add a distinct tested claim to this sequence. A centrality/confidence
association is also omitted: self-trust enters the influence weights, making
that association partly structural rather than an independent mechanism test.

\textbf{Reporting precision.} Rates in the text are rounded summaries, not
confidence intervals. The complexity and temporal analyses use different
weighting, as listed in Table~\ref{tab:samples}. The organizational threshold
range comes from the retained cached grid (positive: 0.673--0.714;
negative: 0.710--0.910). Its target uses surrogate-predicted stationary fractions;
the parity check instead concerns final recorded fractions. Confirm replicate
handling, surrogate approximation and threshold estimation in the methods
before submission. Additional simulation
cases do not substitute for independent tracks, networks or empirical social
observations.
\fi

\end{document}
